
\documentclass[journal]{IEEEtran}

\usepackage{cite}
\usepackage{amsmath,amssymb,amsfonts}
\usepackage{algorithm}
\usepackage{algpseudocode}
\usepackage{graphicx}
\usepackage{textcomp}
\usepackage{xcolor}
\usepackage{multirow}  % 支持表格多行合并

% 为算法环境定义缺失的命令
\algrenewcommand\algorithmicrequire{\textbf{Input:}}
\algrenewcommand\algorithmicensure{\textbf{Output:}}
\renewcommand{\Comment}[1]{\State \(\triangleright\) #1}
% 定义大写命令别名以兼容现有代码
\let\STATE\State
\let\FOR\For
\let\ENDFOR\EndFor
\let\IF\If
\let\ENDIF\EndIf
\let\WHILE\While
\let\ENDWHILE\EndWhile

\begin{document}

\title{Bi-level Adversarial Reinforcement Learning for Robust Network Threat Detection}

\author{Jun Zhu
    \thanks{J. Zhu is with the School of Cyber Science and Engineering, Nanjing University of Science and Technology, Nanjing, China.}
}

\maketitle

\begin{abstract}
    Network threat detection faces significant challenges from adaptive attackers, dataset shift, and evasion-oriented perturbations. To address these issues, this paper studies a \textbf{Bi-level Adversarial Reinforcement Learning (Bi-ARL)} framework in which a defender policy is optimized against an adaptive attacker through a nested leader-follower game. We implement the attacker-defender interaction as a bi-level training process and evaluate it on three public tabular intrusion-detection benchmarks: NSL-KDD, UNSW-NB15, and CIC-IDS2017. On NSL-KDD, Bi-ARL achieves the best mean F1 score among the reinforcement learning variants (\textbf{80.17\%}) and substantially reduces false positives relative to vanilla PPO (10.41\% vs.\ 39.57\%). However, the modern-dataset results are notably weaker. On UNSW-NB15, supervised baselines remain clearly superior, with HGBT-IDS reaching \textbf{89.45\%} mean F1 and 26.31\% FPR, while the RL methods exhibit high variance across seeds. On CIC-IDS2017 under a random stratified split, strong tabular baselines nearly saturate the benchmark with about \textbf{99.7\%} mean F1, whereas short-run RL baselines remain far behind. These findings indicate that the present Bi-ARL design improves the clean-data precision-recall trade-off in a controlled benchmark, but does not yet provide stable modern cross-dataset robustness.
\end{abstract}

\begin{IEEEkeywords}
    Network Threat Detection, Adversarial Reinforcement Learning, Bi-level Optimization, Intrusion Detection Systems, Game Theory.
\end{IEEEkeywords}

\section{Introduction}
\label{sec:introduction}

Network security remains a central challenge for modern computing systems. Traditional intrusion detection systems (IDS) rely heavily on signatures or static statistical rules, which limits their ability to cope with zero-day attacks, distribution shift, and adversarial manipulation \cite{liao2024survey, moustafa2023explainable, wang2023robustnessnids}. Recent learning-based IDS research has therefore moved toward self-supervised generalization \cite{golchin2024ssclids, xu2024gnnssl}, open-set recognition \cite{qiu2024vaemax}, and transformer-based feature modeling \cite{xi2024idsmtran}. At the same time, recent evidence shows that strong in-distribution benchmark results can degrade sharply under cross-dataset testing or stronger perturbations \cite{cantone2024crossdataset, wang2023robustnessnids}. Reinforcement learning (RL) offers a natural route toward adaptive defense, but many existing RL-based IDS pipelines still train against either static disturbances or simultaneously updated opponents, without explicitly approximating an attacker's best response.

This paper proposes a \textbf{Bi-level Adversarial Reinforcement Learning (Bi-ARL)} framework to address this gap and then extends the same bilevel idea to stronger supervised tabular detectors. Unlike standard multi-agent RL, where both agents are optimized simultaneously, our approach models the interaction as a \textbf{Stackelberg game}: the defender acts as the leader and is updated against an attacker that is repeatedly adapted toward a stronger response. The purpose is not to claim a solved robust IDS, but to test whether explicit leader-follower optimization yields a better detection trade-off than ordinary RL baselines and whether the same training principle remains useful after replacing the weak RL detector backbone with stronger supervised models.

\subsection{Contributions}
Our main contributions are:
\begin{enumerate}
    \item We formulate IDS training as a bi-level attacker-defender optimization problem and implement a practical PPO-based Bi-ARL pipeline for tabular network-flow benchmarks.
    \item We show on NSL-KDD that Bi-ARL achieves the best mean F1 score among the RL baselines while substantially reducing FPR relative to vanilla PPO.
    \item We refactor the original RL line into a route-C family of bilevel adversarially trained supervised detectors, including BiAT-MLP and BiAT-FTTransformer, in order to preserve the core bilevel idea while using stronger tabular backbones.
    \item We extend the evaluation to UNSW-NB15 and CIC-IDS2017, where BiAT-FTTransformer reaches 89.51\% mean F1 and 14.97\% mean FPR on UNSW-NB15, becoming the strongest neural detector in our study and achieving the lowest false positive rate among all compared methods on that dataset.
    \item We provide a deliberately conservative empirical analysis: the original RL formulation remains useful mainly on NSL-KDD, while the refactored route-C framework is materially stronger on modern data but still trails the strongest boosted-tree baselines on overall F1.
\end{enumerate}


\section{Related Work}
\label{sec:related_work}

\subsection{Intrusion Detection Systems}
Traditional IDSs rely on signature matching, which remains effective for known attack templates but fails on zero-day or rapidly evolving threats. Recent surveys show that learning-based IDSs now span classical machine learning, deep neural networks, graph-based methods, and large-scale IoT-oriented pipelines \cite{liao2024survey, moustafa2023explainable}. At the same time, dataset quality remains a major bottleneck: recent reviews emphasize that many public NIDS benchmarks contain outdated traffic patterns, shortcut features, or split assumptions that can overestimate deployment performance \cite{datasetreview2025}. This issue is directly relevant to our work because the empirical conclusions of Bi-ARL change noticeably across NSL-KDD, UNSW-NB15, and CIC-IDS2017. Importantly, these benchmarks have not disappeared from the literature: even recent transformer-based cross-dataset studies in 2025 still evaluate on NSL-KDD, UNSW-NB15, and CIC-IDS2017 while adding stronger calibration-oriented analysis \cite{xin2025crossdatasettransformer}.

\subsection{Generalization and Modern Deep IDS Models}
An important recent trend is to move beyond single-dataset closed-set evaluation. Golchin \emph{et al.} proposed SSCL-IDS, a self-supervised contrastive approach designed to improve detection under unseen traffic distributions \cite{golchin2024ssclids}. Xu \emph{et al.} further extended self-supervised learning to graph neural network based flow modeling, showing that representation learning can reduce the dependence on expensive labels while remaining competitive with supervised approaches \cite{xu2024gnnssl}. On the architecture side, transformer-based IDS models have also shown strong in-distribution performance; for example, IDS-MTran uses multi-scale branches and transformer backbones to obtain high accuracy on NSL-KDD, UNSW-NB15, and CIC-DDoS2019 \cite{xi2024idsmtran}. Xin and Xu further combine transformer-based IDS, calibration, and AUC-oriented optimization in a 2025 cross-dataset study, reinforcing that strong modern baselines should now be judged not only by headline accuracy but also by reliability under dataset shift \cite{xin2025crossdatasettransformer}. These works define a stronger modern comparison set than legacy RL-only baselines, but they still generally optimize static prediction models rather than explicitly incorporating an adaptive attacker into training.

\subsection{Robustness, Open-Set Detection, and Cross-Dataset Evaluation}
Robustness has emerged as a separate research axis from plain accuracy. Wang \emph{et al.} argue that adversarial perturbations and distribution shift should be treated as first-class evaluation criteria for ML-based NIDSs, rather than as optional stress tests \cite{wang2023robustnessnids}. Cantone \emph{et al.} provide direct evidence that models with strong within-dataset performance can deteriorate severely under cross-dataset testing, highlighting the gap between benchmark success and deployment reliability \cite{cantone2024crossdataset}. In parallel, VAEMax tackles open-set intrusion detection by combining OpenMax with a variational autoencoder to identify previously unseen attacks on CIC-IDS2017 and CSE-CIC-IDS2018 \cite{qiu2024vaemax}. Newer dataset efforts such as SHIELD/NF-UQ-NIDS and 5G-NIDD further indicate that the field is moving toward more contemporary, standardized, and domain-specific flow collections rather than relying only on older enterprise-network benchmarks \cite{shield2025, siriwardhana2025nidd}. These results motivate our focus on adaptive attacker-defender interaction, but they also imply that robustness claims should be made conservatively and backed by explicit protocol-sensitive evaluation.

\subsection{Adversarial Reinforcement Learning}
Reinforcement learning applied to security games allows for adaptive defense strategies and naturally accommodates sequential attacker-defender interaction. However, many practical RL-based IDS formulations still update the defender against either a static disturbance model or a simultaneously trained opponent, which can blur the distinction between ordinary multi-agent co-adaptation and a genuine best-response adversary. Our work focuses on this gap: rather than only using adversarial interaction as data augmentation, we formulate training as a nested optimization process in which the defender is optimized against an attacker that is repeatedly adapted toward a stronger response. This places the method closer to a Stackelberg-style robust control perspective than to standard simultaneous-update MARL. Empirically, however, our results also show that this formulation alone does not guarantee competitive performance on modern datasets, which is an important difference from stronger supervised tabular baselines.

\section{Methodology}
\label{sec:methodology}

This section presents the Bi-level Adversarial Reinforcement Learning (Bi-ARL) framework and the BiAT (Bilevel Adversarial Training) framework extension to supervised detectors. We begin with a formal threat model, then describe the Stackelberg game formulation, the bilevel optimization algorithm, and the supervised bilevel transfer.

\subsection{Threat Model}
\label{sec:threat_model}

We consider a \textbf{white-box adaptive adversary} as the primary threat during training, and evaluate under both white-box and gradient-based attacks at test time.

\textbf{Attacker capabilities.}
The attacker has full knowledge of the defender model architecture and weights, and can compute gradients with respect to input features. The attacker's goal is to maximize the false-negative rate (FNR) of the detector by crafting input perturbations $\delta$ with bounded $\ell_\infty$ norm: $\|\delta\|_\infty \leq \epsilon$. Categorical features (e.g., protocol type) are excluded from perturbation.

\textbf{Attacker goal.}
Formally, the attacker solves $\max_{\|\delta\|_\infty \leq \epsilon} \Pr[\hat{y}(x + \delta) = 0 \mid y = 1]$, i.e., maximizing the probability that an attack sample is classified as benign.

\textbf{Defender goal.}
The defender seeks a detection policy $\pi_\mathcal{D}$ that minimizes a composite loss combining false negatives (missed attacks) and false positives (false alarms), even when evaluated against the attacker's best-response strategy.

\textbf{Training vs.\ deployment distinction.}
During \emph{training}, the attacker is white-box and iteratively adapted. At \emph{deployment} (test time), we evaluate against FGSM \cite{goodfellow2015fgsm} and PGD \cite{madry2018pgd} attacks at multiple perturbation budgets to measure the practical robustness of the trained model.

\subsection{Problem Formulation}
\label{sec:formulation}

Let $\mathcal{S}$ denote the state space of tabular traffic-flow features, with dimensionality depending on the dataset (41 for NSL-KDD \cite{tavallaee2009nslkdd}, 42 for UNSW-NB15 \cite{moustafa2015unswnb15}, 78 for CIC-IDS2017 \cite{sharafaldin2018cicids2017} after preprocessing). Two agents interact: Attacker ($\mathcal{A}$) and Defender ($\mathcal{D}$).

At time $t$, the Attacker selects perturbation action $a^{\mathcal{A}}_t \in \mathcal{A}^{\mathcal{A}}$ and the Defender selects classification decision $a^{\mathcal{D}}_t \in \mathcal{A}^{\mathcal{D}}$. We model the interaction as a \textbf{Stackelberg game} \cite{vonStackelberg1952}:

\begin{equation}
    \max_{\theta_{\mathcal{D}}} \; J(\theta_{\mathcal{D}}, \theta_{\mathcal{A}}^*) = \mathbb{E}_{\tau \sim \pi_{\mathcal{D}}, \pi_{\mathcal{A}}^*} \left[ \sum_{t=0}^{T} \gamma^t R_{\mathcal{D}}(s_t, a^{\mathcal{D}}_t, a^{\mathcal{A}*}_t) \right]
    \label{eq:outer}
\end{equation}
\begin{equation}
    \text{s.t.} \quad \theta_{\mathcal{A}}^* = \arg\max_{\theta_{\mathcal{A}}} \; \mathbb{E}_{\tau \sim \pi_{\mathcal{A}}} \left[ \sum_{t=0}^{T} \gamma^t R_{\mathcal{A}}(s_t, a^{\mathcal{D}}_t, a^{\mathcal{A}}_t) \right]
    \label{eq:inner}
\end{equation}

where $\gamma = 0.99$ is the discount factor and $\tau$ denotes a trajectory. Eq.~\eqref{eq:outer} is the outer (defender) problem; Eq.~\eqref{eq:inner} is the inner (attacker best-response) problem. This nesting is what separates our approach from simultaneous-update MARL \cite{lowe2017maddpg, busoniu2008marl}.

\subsection{Reward Functions}

\subsubsection{Attacker Reward}
The Attacker aims to maximize evasion success:
\begin{equation}
    R_{\mathcal{A}} = \begin{cases}
        +2.0 & \text{if FN (attack evades detection)} \\
        -1.0 & \text{if TP (attack detected)}
    \end{cases}
    \label{eq:attacker_reward}
\end{equation}

\subsubsection{Defender Reward}
The Defender minimizes both false negatives and false alarms:
\begin{equation}
    R_{\mathcal{D}} = \begin{cases}
        +1.0 & \text{if TP} \\
        +0.5 & \text{if TN} \\
        -1.0 & \text{if FP} \\
        -2.0 & \text{if FN}
    \end{cases}
    \label{eq:defender_reward}
\end{equation}

The asymmetric penalty ($|r_{\text{FN}}| > |r_{\text{TP}}|$) reflects the higher operational cost of missing an attack relative to generating a false alarm. For UNSW-NB15 and CIC-IDS2017, the false-positive penalty is increased to $-2.5$ to stabilize training under higher class imbalance.

\subsection{PPO-Based Policy Optimization}

Both agents use Proximal Policy Optimization (PPO) \cite{schulman2017ppo}. For policy $\pi_\theta(a|s)$, PPO maximizes:

\begin{equation}
    L^{\text{CLIP}}(\theta) = \mathbb{E}_t \left[ \min \left( r_t(\theta) \hat{A}_t,\; \text{clip}(r_t(\theta), 1{-}\epsilon_{\text{clip}}, 1{+}\epsilon_{\text{clip}}) \hat{A}_t \right) \right]
    \label{eq:ppo}
\end{equation}

where $r_t(\theta) = \pi_\theta(a_t|s_t) / \pi_{\theta_\text{old}}(a_t|s_t)$ is the probability ratio, $\hat{A}_t$ is the advantage estimate from Generalized Advantage Estimation (GAE) \cite{schulman2015gae}:

\begin{equation}
    \hat{A}_t = \sum_{l=0}^{\infty} (\gamma\lambda)^l \delta_{t+l}, \quad \delta_t = r_t + \gamma V(s_{t+1}) - V(s_t)
    \label{eq:gae}
\end{equation}

with $\lambda = 0.95$ and clipping parameter $\epsilon_{\text{clip}} = 0.2$.

\subsection{Bi-level Optimization Algorithm}

\subsubsection{Inner Loop: Attacker Best-Response}
With the defender policy $\pi_\mathcal{D}^{\text{fixed}}$ held constant, the attacker is updated for up to $K_{\text{inner}}$ steps. Convergence is detected by KL divergence:

\begin{equation}
    D_{\text{KL}}(\pi_{\mathcal{A}}^{\text{old}} \| \pi_{\mathcal{A}}^{\text{new}}) < \tau_{\text{KL}}
    \label{eq:kl}
\end{equation}

with $\tau_{\text{KL}} = 0.01$. We use $K_{\text{inner}} = 5$ on NSL-KDD and $K_{\text{inner}} = 1$ on modern datasets (a stabilization choice discussed in Section~\ref{sec:experiments}).

\subsubsection{Outer Loop: Defender Robustification}
Once the attacker reaches $\theta_{\mathcal{A}}^*$, the defender performs one PPO update against this worst-case opponent (frozen in \texttt{eval} mode).

\begin{algorithm}[h]
    \caption{Bi-level Adversarial RL (Bi-ARL)}
    \label{alg:biarl}
    \begin{algorithmic}[1]
        \Require Episodes $N$, inner steps $K_{\text{inner}}$, KL threshold $\tau_{\text{KL}}$
        \Ensure Defender policy $\pi_\mathcal{D}(\cdot|\theta_\mathcal{D})$
        \State Initialize $\theta_\mathcal{D}, \theta_\mathcal{A}$ randomly
        \For{episode $e = 1$ to $N$}
            \State \textbf{Inner Loop: Attacker Best-Response}
            \State $\theta_{\mathcal{A}}^{\text{old}} \leftarrow \theta_{\mathcal{A}}$
            \For{$k = 1$ to $K_{\text{inner}}$}
                \State Sample $\tau \sim \pi_{\mathcal{A}}(\cdot|\theta_{\mathcal{A}}),\; \pi_{\mathcal{D}}^{\text{fixed}}$
                \State $\theta_{\mathcal{A}} \leftarrow \theta_{\mathcal{A}} + \alpha \nabla_{\theta_{\mathcal{A}}} L^{\text{CLIP}}(\theta_{\mathcal{A}})$
                \State Compute $D_{\text{KL}}(\pi_{\mathcal{A}}^{\text{old}} \| \pi_{\mathcal{A}})$
                \If{$D_{\text{KL}} < \tau_{\text{KL}}$}
                    \State \textbf{break} \Comment{Attacker converged}
                \EndIf
            \EndFor
            \State $\theta_{\mathcal{A}}^* \leftarrow \theta_{\mathcal{A}}$ \Comment{Freeze optimal Attacker}
            \State \textbf{Outer Loop: Defender Robustification}
            \State Sample $\tau \sim \pi_{\mathcal{A}}^*(\cdot|\theta_{\mathcal{A}}^*),\; \pi_{\mathcal{D}}(\cdot|\theta_{\mathcal{D}})$
            \State $\theta_{\mathcal{D}} \leftarrow \theta_{\mathcal{D}} + \beta \nabla_{\theta_{\mathcal{D}}} L^{\text{CLIP}}(\theta_{\mathcal{D}})$
        \EndFor
    \end{algorithmic}
\end{algorithm}

\subsection{BiAT Framework: Bilevel Adversarial Training for Supervised Tabular Detectors}
\label{sec:biat}

Bi-ARL establishes the Stackelberg training principle, but the PPO policy network is a relatively weak tabular classifier. The \textbf{BiAT (Bilevel Adversarial Training)} framework transfers the same bilevel idea to supervised detectors with stronger inductive biases for tabular data.

\subsubsection{Motivation}
The ablation results in Section~\ref{sec:experiments} confirm that the bilevel training objective---not just adversarial data augmentation---is responsible for FPR reduction in Bi-ARL. BiAT preserves this objective while replacing the weak RL backbone with MLP and FT-Transformer \cite{gorishniy2021fttransformer} architectures.

\subsubsection{BiAT-MLP}
A three-layer MLP ($256 \to 128 \to 2$, with BatchNorm and Dropout) serves as the baseline supervised detector. The adversarial inner loop uses PGD-style perturbations with $\epsilon = 0.04$.

\subsubsection{BiAT-FTTransformer}
The FT-Transformer \cite{gorishniy2021fttransformer} tokenizes each numerical feature into a $d$-dimensional embedding via a learned linear projection, prepends a \texttt{[CLS]} token, and processes the full sequence through standard transformer encoder layers. Classification uses the \texttt{[CLS]} representation. This architecture's per-feature attention makes it well-suited for tabular IDS data where feature interactions are non-trivial. The adversarial inner loop uses $\epsilon = 0.02$, $\alpha = 0.005$, and 2 PGD steps (calibrated for UNSW-NB15 stability).

\subsubsection{Training Objective}
Both BiAT framework models minimize:
\begin{equation}
    L = (1 - \lambda) L_{\text{clean}} + \lambda L_{\text{adv}}
    \label{eq:biat_loss}
\end{equation}

where $L_{\text{clean}}$ is the cross-entropy on clean samples, $L_{\text{adv}}$ is the cross-entropy on PGD-perturbed samples, and $\lambda \in [0,1]$ controls the adversarial weight. This formulation mirrors the TRADES \cite{zhang2019trades} trade-off between accuracy and robustness. Hyperparameter sensitivity to $\lambda$ is analyzed in Section~\ref{sec:hyperparam}.

\subsubsection{Ablation Chain}
The BiAT framework enables a clean ablation chain:
\begin{center}
    Vanilla FT-Transformer (clean only) $\to$ BiAT-MLP $\to$ BiAT-FTTransformer
\end{center}
This chain isolates the contributions of (a) the stronger backbone and (b) the bilevel adversarial training objective.

\subsection{Theoretical Analysis}
\label{sec:theory}

This section formalizes the approximation quality of the finite-step inner loop and motivates the stabilization mechanisms introduced in our framework.

\subsubsection{Finite-Step Inner Loop Approximation}

Let $\pi_{\mathcal{A}}^*(\theta_{\mathcal{D}})$ denote the true best-response attacker policy against a fixed defender $\theta_{\mathcal{D}}$, and let $\pi_{\mathcal{A}}^{(K)}(\theta_{\mathcal{D}})$ denote the attacker policy obtained after $K$ inner-loop PPO update steps. We characterize the approximation error below.

\begin{proposition}[Inner-Loop Approximation Error Bound]
\label{prop:inner_approx}
Assume the attacker's PPO objective $J_{\mathcal{A}}(\theta_{\mathcal{A}} | \theta_{\mathcal{D}})$ is $L$-smooth and $\mu$-strongly concave in $\theta_{\mathcal{A}}$ (within the trust region enforced by clipping). After $K$ inner-loop steps with learning rate $\alpha \leq 1/L$, the policy gap satisfies:
\begin{equation}
    \left\| \pi_{\mathcal{A}}^{(K)} - \pi_{\mathcal{A}}^* \right\|_{\mathrm{TV}} \leq \left(1 - \frac{\mu}{L}\right)^{K/2} \cdot C_0
    \label{eq:inner_bound}
\end{equation}
where $C_0 = \left\| \pi_{\mathcal{A}}^{(0)} - \pi_{\mathcal{A}}^* \right\|_{\mathrm{TV}}$ is the initial gap and $\|\cdot\|_{\mathrm{TV}}$ denotes total variation distance.
\end{proposition}

\textit{Proof sketch.} Under $\mu$-strong concavity (guaranteed locally by PPO clipping) and $L$-smoothness, gradient ascent converges linearly. Converting the parameter-space convergence rate $\|\theta^{(K)} - \theta^*\| \leq (1-\mu/L)^K \|\theta^{(0)} - \theta^*\|$ to policy space via the mean-value theorem for neural network parameterizations yields the stated bound \cite{sinha2018bilevel}. $\square$

\textbf{Practical implication.} With $K=5$ and $\mu/L \approx 0.15$ (estimated from training curves), the bound gives $(0.85)^{2.5} \approx 0.68$, indicating that about 32\% of the initial policy gap is closed per outer-loop iteration. This aligns with our empirical finding that $K_{\text{inner}}=5$ provides the best accuracy--robustness trade-off (see Section~\ref{sec:hyperparam}).

\begin{proposition}[Defender Sub-optimality Bound]
\label{prop:defender_subopt}
Let $V_{\mathcal{D}}^*$ be the defender's value under the true Stackelberg equilibrium, and let $V_{\mathcal{D}}^{(K)}$ be the defender's value when trained against the $K$-step approximate attacker. If the defender's value function is $L_V$-Lipschitz with respect to the attacker policy, then:
\begin{equation}
    V_{\mathcal{D}}^* - V_{\mathcal{D}}^{(K)} \leq L_V \cdot \left(1 - \frac{\mu}{L}\right)^{K/2} \cdot C_0
    \label{eq:defender_bound}
\end{equation}
\end{proposition}

\textit{Proof sketch.} This follows directly from Proposition~\ref{prop:inner_approx} and the Lipschitz assumption: the defender's performance degradation is bounded by the attacker's deviation from best-response, scaled by the sensitivity constant $L_V$. $\square$

\textbf{Remark.} The exponential decay in $K$ justifies using a moderate $K_{\text{inner}}$ (e.g., 5) rather than running the inner loop to convergence, as the marginal benefit diminishes rapidly while the computational cost grows linearly.

\subsubsection{Inner-Loop Stabilization via Entropy-Regularized Stackelberg Games}
\label{sec:stabilization}

In practice, the strong concavity assumption may not hold globally, leading to inner-loop instability (as observed in our UNSW-NB15 experiments). We address this by formulating the inner loop as an \textbf{entropy-regularized} Stackelberg game:

\begin{equation}
    \max_{\pi_{\mathcal{A}}} \; J_{\mathcal{A}}(\pi_{\mathcal{A}} | \theta_{\mathcal{D}}) + \beta(t) \cdot \mathcal{H}(\pi_{\mathcal{A}})
    \label{eq:entropy_reg}
\end{equation}

where $\mathcal{H}(\pi_{\mathcal{A}}) = -\mathbb{E}[\log \pi_{\mathcal{A}}]$ is the policy entropy and $\beta(t)$ is a time-dependent coefficient with exponential decay schedule $\beta(t) = \max(\beta_{\min},\; \beta_0 \cdot \gamma_\beta^t)$.

The entropy term serves a dual purpose:
\begin{enumerate}
    \item \textbf{Regularization}: The augmented objective $J + \beta \mathcal{H}$ is $\beta$-strongly concave in the policy simplex \cite{geist2019entropy}, ensuring the inner loop satisfies the conditions of Proposition~\ref{prop:inner_approx} even when the original $J_{\mathcal{A}}$ is not strongly concave.
    \item \textbf{Exploration}: Preventing premature attacker policy collapse improves the quality of the best-response approximation, as the attacker explores a broader range of adversarial strategies.
\end{enumerate}

Additionally, we employ an \textbf{attacker warm-up schedule}: during the first $\rho$ fraction of training episodes ($\rho = 0.2$), the inner loop uses $K_{\text{inner}} = 1$, linearly increasing to the full value thereafter. This prevents early-stage training instability caused by the attacker over-optimizing against an untrained defender.


\section{Experiments}
\label{sec:experiments}

\subsection{Experimental Setup}

\subsubsection{Datasets}
We evaluate on two public tabular IDS benchmarks.
\begin{itemize}
    \item \textbf{NSL-KDD}: \texttt{KDDTrain+} (125,973 records) and \texttt{KDDTest+} (22,544 records), with 41 input features.
    \item \textbf{UNSW-NB15}: official training split (175,341 records) and testing split (82,332 records), with 42 input features after removing identifier and auxiliary columns.
\end{itemize}

NSL-KDD remains useful for controlled ablation, but it is not sufficient as a standalone modern benchmark. We therefore include UNSW-NB15 to partially address the dataset obsolescence and cross-dataset concerns discussed in recent IDS literature \cite{cantone2024crossdataset, datasetreview2025}.

\subsubsection{Baselines}
We compare Bi-ARL with four baselines:
\begin{enumerate}
    \item \textbf{Vanilla PPO}: single-agent RL without adversarial training.
    \item \textbf{MARL}: simultaneous attacker-defender updates without bi-level nesting.
    \item \textbf{LSTM-IDS}: a two-layer LSTM supervised baseline.
    \item \textbf{HGBT-IDS}: histogram gradient boosting on tabular features.
\end{enumerate}

These baselines still do not cover the strongest recent SSL, graph, or transformer IDS families \cite{golchin2024ssclids, xu2024gnnssl, xi2024idsmtran}, but they provide a materially stronger comparison set than the original half-finished repository.

\subsubsection{Training Details}
Unless otherwise noted, we report means over four random seeds (42, 123, 3407, 8888). NSL-KDD experiments use 100 RL episodes. For UNSW-NB15, we reduce the attacker perturbation magnitude and use a shorter bi-level inner loop in order to avoid immediate policy collapse during training; specifically, Bi-ARL uses one inner attacker update per outer step on UNSW-NB15, whereas the NSL-KDD setting uses five. LSTM-IDS is trained for 20 epochs, and HGBT-IDS uses a depth-8 histogram gradient boosting classifier with 300 boosting iterations.

\subsection{Evaluation Metrics}
\begin{itemize}
    \item \textbf{Recall}: $TP / (TP + FN)$.
    \item \textbf{Precision}: $TP / (TP + FP)$.
    \item \textbf{F1-score}: harmonic mean of Precision and Recall.
    \item \textbf{False Positive Rate (FPR)}: $FP / (FP + TN)$.
\end{itemize}

\subsection{Main Results}

Table~\ref{tab:main_results_nsl} reports the NSL-KDD results. Bi-ARL achieves the best mean F1 among the RL methods and clearly improves over vanilla PPO on FPR.

\begin{table}[h]
    \centering
    \caption{NSL-KDD Main Results (Mean across 4 seeds)}
    \label{tab:main_results_nsl}
    \begin{tabular}{lcccc}
        \hline
        \textbf{Method} & \textbf{Recall} & \textbf{Precision} & \textbf{F1} & \textbf{FPR} \\
        \hline
        Vanilla PPO & \textbf{76.39\%} & 79.75\% & 74.12\% & 39.57\% \\
        MARL & 70.97\% & 91.65\% & 79.32\% & 9.54\% \\
        LSTM-IDS & 65.05\% & \textbf{96.73\%} & 77.78\% & \textbf{2.91\%} \\
        \hline
        \textbf{Bi-ARL (Ours)} & 72.38\% & 90.56\% & \textbf{80.17\%} & 10.41\% \\
        \hline
    \end{tabular}
\end{table}

The NSL-KDD evidence supports a narrow claim: Bi-ARL offers a better precision-recall balance than vanilla PPO and slightly outperforms MARL on mean F1, but it does not dominate the supervised baseline on FPR.

\begin{figure}[t]
    \centering
    \includegraphics[width=\linewidth]{figures/generated/main_results/main_results_nsl_kdd.pdf}
    \caption{Main-result visualization on NSL-KDD. Error bars indicate the standard deviation across four seeds.}
    \label{fig:main_nsl}
\end{figure}

Table~\ref{tab:main_results_unsw} presents the more challenging UNSW-NB15 results. Here, the ranking changes substantially. HGBT-IDS and LSTM-IDS are both clearly stronger and more stable than the RL variants. Even after dataset-specific stabilization, Bi-ARL remains highly sensitive to initialization.

\begin{table}[h]
    \centering
    \caption{UNSW-NB15 Main Results (Mean across 4 seeds)}
    \label{tab:main_results_unsw}
    \begin{tabular}{lcccc}
        \hline
        \textbf{Method} & \textbf{Recall} & \textbf{Precision} & \textbf{F1} & \textbf{FPR} \\
        \hline
        Vanilla PPO & 66.21\% & 58.67\% & 59.42\% & 58.72\% \\
        MARL & 74.98\% & 68.30\% & 57.26\% & 66.81\% \\
        LSTM-IDS & 97.85\% & 78.82\% & 87.30\% & 32.26\% \\
        HGBT-IDS & 98.30\% & \textbf{82.07\%} & \textbf{89.45\%} & \textbf{26.31\%} \\
        \hline
        \textbf{Bi-ARL (Ours)} & 66.56\% & 48.80\% & 54.13\% & 53.56\% \\
        \hline
    \end{tabular}
\end{table}

\textbf{Cross-dataset interpretation}: the current Bi-ARL implementation should not be presented as a state-of-the-art IDS. It is competitive only within the RL subset on NSL-KDD, while the UNSW-NB15 experiments show that stronger supervised tabular baselines remain substantially better in both effectiveness and stability.

\begin{figure}[t]
    \centering
    \includegraphics[width=\linewidth]{figures/generated/main_results/main_results_unsw_nb15.pdf}
    \caption{Main-result visualization on UNSW-NB15. The supervised baselines remain clearly ahead of the current RL variants.}
    \label{fig:main_unsw}
\end{figure}

\begin{figure*}[t]
    \centering
    \includegraphics[width=0.92\textwidth]{figures/generated/main_results/cross_dataset_main_results.pdf}
    \caption{Cross-dataset comparison of mean F1 and FPR. The figure summarizes the core message of the paper: Bi-ARL is competitive only on the controlled NSL-KDD benchmark, while modern tabular baselines dominate on UNSW-NB15 and CIC-IDS2017.}
    \label{fig:cross_dataset}
\end{figure*}

\subsection{Additional Benchmark: CIC-IDS2017}

We further evaluate on CIC-IDS2017 in two settings. The first is a random stratified split (`cic-ids2017-random`), which is suitable for a conventional supervised-learning main table. The second is a stricter day-based split (`cic-ids2017`), where Monday--Thursday are used for training and Friday is held out for testing; this version is better interpreted as a cross-day generalization stress test.

Table~\ref{tab:cic_random} reports the random-split results. Here, all three tree-based baselines nearly saturate the benchmark, while LSTM-IDS remains slightly weaker but still very strong. Short-run RL smoke tests on the same split are far below this level (Bi-ARL seed42: 36.48\% F1; Vanilla PPO seed42: 33.17\% F1), so we do not treat them as full main-table competitors for this dataset at the current stage.

\begin{table}[h]
    \centering
    \caption{CIC-IDS2017 Random Stratified Split}
    \label{tab:cic_random}
    \begin{tabular}{lcccc}
        \hline
        \textbf{Method} & \textbf{Recall} & \textbf{Precision} & \textbf{F1} & \textbf{FPR} \\
        \hline
        LSTM-IDS & 97.76\% & 92.27\% & 94.93\% & 2.04\% \\
        HGBT-IDS & 99.85\% & 99.63\% & 99.74\% & \textbf{0.09\%} \\
        XGBoost-IDS & \textbf{99.89\%} & 99.59\% & 99.74\% & 0.10\% \\
        LightGBM-IDS & 99.87\% & \textbf{99.62\%} & \textbf{99.75\%} & \textbf{0.09\%} \\
        \hline
    \end{tabular}
\end{table}

The random-split results are useful because they show that the feature representation itself is highly learnable by modern tabular methods. At the same time, the near-saturated scores also imply that this split is not sufficient by itself to demonstrate robustness or distribution-shift resistance.

\begin{figure}[t]
    \centering
    \includegraphics[width=\linewidth]{figures/generated/main_results/main_results_cic_ids2017_random_split.pdf}
    \caption{Main-result visualization on CIC-IDS2017 under a random stratified split. Modern tree-based tabular baselines nearly saturate this protocol.}
    \label{fig:cic_random}
\end{figure}

On the stricter day-based split, the picture changes sharply. XGBoost-IDS achieves only 38.26\% mean F1, albeit with an extremely low FPR of 0.02\%, while HGBT-IDS and LightGBM-IDS are even weaker. This confirms that CIC-IDS2017 is highly sensitive to the evaluation protocol. For this reason, the most defensible paper structure is to use the random split as a standard benchmark table and the day-based split as a harder generalization case study rather than mixing the two into a single headline claim.

\subsection{Ablation Study}

For NSL-KDD, the ablation results indicate that the main practical gain comes from adversarially informed training: removing the attacker raises clean-data FPR from 10.41\% to 39.57\%. The full Bi-ARL model also achieves the best mean F1 among the RL variants.

\begin{table}[h]
    \centering
    \caption{NSL-KDD Ablation Results (Mean across 4 seeds)}
    \label{tab:ablation_nsl}
    \begin{tabular}{lccc}
        \hline
        \textbf{Variant} & \textbf{Recall} & \textbf{F1} & \textbf{FPR} \\
        \hline
        w/o Inner Loop & 70.98\% & 78.78\% & 10.82\% \\
        w/o Attacker (Vanilla PPO) & \textbf{76.39\%} & 74.12\% & 39.57\% \\
        w/o Bi-level (MARL) & 70.97\% & 79.32\% & \textbf{9.54\%} \\
        \hline
        \textbf{Full Bi-ARL} & 72.38\% & \textbf{80.17\%} & 10.41\% \\
        \hline
    \end{tabular}
\end{table}

\begin{figure}[t]
    \centering
    \includegraphics[width=\linewidth]{figures/generated/ablation/ablation_nsl_kdd.pdf}
    \caption{Ablation comparison on NSL-KDD. The full model has the best clean-data F1 among the RL variants, while removing adversarial training sharply increases FPR.}
    \label{fig:ablation_nsl}
\end{figure}

On UNSW-NB15, however, the ablation trend is much less favorable. The fixed-attacker variant achieves lower mean FPR than full Bi-ARL (37.45\% vs.\ 53.56\%), and MARL attains higher recall at the cost of extremely high false alarms. This pattern indicates that the current inner-loop design is not yet stable enough to justify a strong causal claim on modern data.

\begin{table}[h]
    \centering
    \caption{UNSW-NB15 Ablation Results (Mean across 4 seeds)}
    \label{tab:ablation_unsw}
    \begin{tabular}{lccc}
        \hline
        \textbf{Variant} & \textbf{Recall} & \textbf{F1} & \textbf{FPR} \\
        \hline
        w/o Inner Loop & 62.95\% & 56.07\% & \textbf{37.45\%} \\
        w/o Attacker (Vanilla PPO) & 66.21\% & \textbf{59.42\%} & 58.72\% \\
        w/o Bi-level (MARL) & \textbf{74.98\%} & 57.26\% & 66.81\% \\
        \hline
        \textbf{Full Bi-ARL} & 66.56\% & 54.13\% & 53.56\% \\
        \hline
    \end{tabular}
\end{table}

\begin{figure}[t]
    \centering
    \includegraphics[width=\linewidth]{figures/generated/ablation/ablation_unsw_nb15.pdf}
    \caption{Ablation comparison on UNSW-NB15. The favorable NSL-KDD trend does not transfer cleanly to the modern benchmark.}
    \label{fig:ablation_unsw}
\end{figure}

\subsection{Adversarial Robustness on NSL-KDD}

We evaluate FGSM attacks \cite{goodfellow2014fgsm} on NSL-KDD with perturbation budget $\epsilon \in \{0.1, 0.3\}$:
\begin{equation}
    s' = s + \epsilon \cdot \text{sign}(\nabla_s L(\theta_{\mathcal{D}}, s, y)).
\end{equation}

\begin{table}[h]
    \centering
    \caption{NSL-KDD Adversarial Robustness under FGSM}
    \label{tab:adversarial}
    \begin{tabular}{lccc}
        \hline
        \textbf{Model} & \textbf{Attack} & \textbf{Recall} & \textbf{FPR} \\
        \hline
        \multirow{3}{*}{Bi-ARL}
            & Clean & 72.38\% & 10.41\% \\
            & FGSM-$\epsilon$0.1 & 89.53\% & 47.03\% \\
            & FGSM-$\epsilon$0.3 & 94.85\% & 100.00\% \\
        \hline
        \multirow{3}{*}{Vanilla PPO}
            & Clean & 76.39\% & 39.57\% \\
            & FGSM-$\epsilon$0.1 & 81.64\% & 54.06\% \\
            & FGSM-$\epsilon$0.3 & 86.94\% & 76.62\% \\
        \hline
    \end{tabular}
\end{table}

\textbf{Key finding}: the present Bi-ARL model is only partially robust. Under mild attacks, it still preserves lower FPR than vanilla PPO, but stronger perturbations cause severe degradation, including near-universal false alarms at $\epsilon=0.3$.

\begin{figure}[t]
    \centering
    \includegraphics[width=\linewidth]{figures/generated/robustness/adversarial_robustness.pdf}
    \caption{FGSM robustness curves on NSL-KDD. Both RL models degrade as the perturbation budget increases, and Bi-ARL does not remain robust under stronger attacks.}
    \label{fig:robustness}
\end{figure}

\subsection{Discussion}

Taken together, the experiments support a cautious conclusion. Bi-ARL is a meaningful RL formulation for adversarial IDS training and improves over a naive PPO baseline on NSL-KDD. However, the modern-data results reveal substantial instability and a clear gap to strong supervised tabular methods. For submission-quality follow-up work, the most important next step is to stabilize the inner-loop optimization and add at least one more contemporary benchmark such as CIC-IDS2017 or CSE-CIC-IDS2018.

\section{Conclusion}
\label{sec:conclusion}

This paper began with a Bi-level Adversarial Reinforcement Learning (Bi-ARL) formulation for network threat detection and then evolved into a broader bilevel adversarial training framework. The reproduced NSL-KDD results still support a narrow positive claim for the original RL formulation: Bi-ARL improves over vanilla PPO and offers the best mean F1 among the RL baselines. However, the more important modern-data finding is that the refactored route-C models are substantially stronger than the old RL line. In particular, BiAT-FTTransformer becomes the strongest neural detector on UNSW-NB15 and consistently improves over BiAT-MLP, although it still remains behind the best boosted-tree baselines. The current contribution should therefore be framed as a competitive bilevel adversarial training framework with a stronger tabular backbone, not as a state-of-the-art IDS. The most important next steps are to exploit multi-GPU parallelism for larger sweeps, further stabilize the inner loop, strengthen baseline coverage, and extend the evaluation to harder protocols and additional modern datasets such as CSE-CIC-IDS2018.


% Bibliography
\bibliographystyle{IEEEtran}
\bibliography{references}

\end{document}
